\subsection{Data provenance, ethics and disclosure}
\label{sec:ethics}

All three datasets are public and need no credentialed data use agreement. Synthea is synthetic,
MIMIC-IV and eICU-CRD de-identified demonstration releases. All three are pinned locally by SHA-256,
so no access request can fail mid-study. No IRB review was sought. The same controls on identifiable
records would need IRB approval and a data use agreement before a single run. That setting is where
the privacy claim would be tested, and this study never enters it.

No re-identification was attempted. The model instrument receives only the projected slice
\texttt{project()} builds from the manifest (Section~\ref{sec:designs}), so no field outside a
manifest reaches it, and that dispatch is the control measured.

Stand-in identifiers come from a seeded counter, never derived from the identifier they replace,
as 45 C.F.R. \S 164.514(c)(1) requires of a re-identification code. Hashing the medical record
number, which most would reach for, would not qualify. That covers one field. It does not cover
the pipeline around it, and the result above is reduced exposure, not de-identification.

AI tooling use is disclosed in \texttt{AI\_DISCLOSURE.md}, and the register of instrument defects,
the items left open and the test-suite critique are in \texttt{codebase\_audit/}.
